\documentclass[UTF8,a4paper,11pt,fontset=none]{ctexart}
\usepackage{fontspec}
\setCJKmainfont[AutoFakeBold=2,AutoFakeSlant=0.2]{SimSun}
\setCJKsansfont{SimHei}
\setCJKmonofont{SimSun}
\setmonofont{Consolas}
\usepackage[left=27mm,right=27mm,top=25mm,bottom=25mm]{geometry}
\usepackage{amsmath,amssymb,amsthm,mathtools}
\usepackage{booktabs,enumitem,microtype,hyperref}
\hypersetup{hidelinks,pdfauthor={宋敏, 宋瑷成, 宋奥菲, 范钧睿, 胡安, 涂源},pdftitle={禁用子集和逆定理的线性范围改进}}
\pagestyle{plain}
\setlength{\parindent}{2em}
\setlength{\parskip}{0.12em}
\linespread{1.12}
\allowdisplaybreaks[2]
\newtheorem{theorem}{定理}[section]
\newtheorem{lemma}[theorem]{引理}
\newtheorem{corollary}[theorem]{推论}
\newtheorem{proposition}[theorem]{命题}
\theoremstyle{remark}
\newtheorem{remark}[theorem]{注}
\newcommand{\N}{\mathbb N}
\newcommand{\Z}{\mathbb Z}
\newcommand{\sig}{\Sigma}
\newcommand{\snd}{\operatorname{snd}}
\newcommand{\gcdop}{\operatorname{gcd}}
\newcommand{\rthree}{r_3}
\newcommand{\eps}{\varepsilon}
\title{\bfseries 禁用子集和逆定理的线性范围改进\\[0.3em]
\large 及 $q=4,\ c=2/3$ 的同余类极限}
\author{宋敏、宋瑷成、宋奥菲、范钧睿、胡安、涂源}
\date{2026 年 11 月 19 日\\[0.3em]\small v0.1.1：署名更新；数学内容沿用 v0.1.0}
\begin{document}
\maketitle
\begin{abstract}
设 $s$ 为正整数，$q=\snd(s)$ 为不整除 $s$ 的最小正整数。
本文将 Matijasevick 的禁用子集和逆定理中的范围因子
$2^{q-2}$ 改为 $q-1$，并相应缩短表示长度。
具体而言，对每个 $\eps>0$，存在仅依赖于 $s,\eps$ 的阈值 $T_0$，
使得 $T\ge T_0$、$(q-1)E\le sT$ 及
$A\subseteq[1,E]$、$|A|\ge E/q+\eps T$ 蕴含：
$sT$ 可表示为至多 $(q-1)(\lceil2s/\eps\rceil+1)$
个互异 $A$ 元素之和。
证明以 Lev 的多重和集增量界替换原论证的二进迭代，
其余使用既有的好中心抽取与重复项消去框架。
由此得到范围常数 $\theta_q\le q-1$；
结合已有模数构造，沿 $n\equiv6\pmod{12}$ 有
$F(\lfloor2n/3\rfloor,n)/n\to1/4$。
本文不将 Lev 路线或所复用的组合桥本身作为新结果。
附录说明有限检查及去重桥的 Lean 证明范围；
主定理的论证为引用 Lev 与 Roth 定理的纸面证明。
\par\noindent\textbf{关键词：}子集和；逆加法定理；多重和集；最小非因子；Roth 定理。
\end{abstract}

\section{引言与既有结果的边界}\label{sec:intro}
对有限正整数集 $A$，定义非空互异元素子集和
\[
  \sig(A)=
  \left\{\sum_{a\in X}a:\ \varnothing\ne X\subseteq A\right\}.
\]
对 $M\ge0$、$n\ge1$，记
\[
 F(M,n)=\max\bigl\{|A|:\ A\subseteq[1,M],\ n\notin\sig(A)\bigr\},
 \qquad f_c(n)=F(\lfloor cn\rfloor,n).
\]
这里及下文 $[u,v]$ 表示整数区间。
令 $\snd(s)$ 表示最小的不整除正整数 $s$ 的正整数。
该数 $q=\snd(s)$ 至少为 $2$，并且每个 $1\le d<q$ 都整除 $s$。

Matijasevick~\cite{Matijasevick2026} 的 Theorem~1.2 表明：
对固定 $s,\eps>0$，当 $T$ 充分大且 $2^{q-2}E\le sT$ 时，
$|A|\ge E/q+\eps T$ 蕴含 $sT\in\sig(A)$；
其给出的表示长度不超过
$2^{q-2}(\lceil2s/\eps\rceil+1)$。
原文定义范围常数 $\theta_q$ 并得到
\[
  \theta_2=1,\quad \theta_3=2,\qquad
  \max\!\left\{2,q-2+\frac1q\right\}\le\theta_q\le2^{q-2}
  \quad(q\ge4),
\]
并询问 $\theta_q$ 是否线性增长（原文 Problem~9.2）。
尤须指出：\textbf{原文 Remark~3.2 已明确建议考察 Lev 型多重和集结构定理}。
因此，本文讨论的具体增量限于 $q-1$ 范围在该禁用子集和问题中的应用
及 $q=4,c=2/3$ 的极限推论，而不是首次提出 Lev 路线。
好中心、去重及中心覆盖的证明是既有组合框架的重述；
有界基结论也可由 Lev~\cite{Lev1997} 的既有区间定理直接推出。
下文保留增量式的短证明，以清楚展示因子 $q-1$ 从何而来。

\begin{theorem}[线性范围逆定理]\label{thm:main}
固定整数 $s\ge1$，令 $q=\snd(s)$、$r=q-1$。
对每个 $\eps>0$，存在整数 $T_0=T_0(s,\eps)$，使得对所有
整数 $T\ge T_0$、$E\ge1$ 与 $A\subseteq[1,E]$，只要
\[
  rE\le sT,\qquad |A|\ge\frac Eq+\eps T,
\]
则 $sT$ 可以表示成 $A$ 中至多
\[
  H=r\bigl(\lceil2s/\eps\rceil+1\bigr)
\]
个互异元素之和。
\end{theorem}

\begin{theorem}[指定同余类上的极限]\label{thm:limit}
沿正整数 $n\equiv6\pmod{12}$ 趋于无穷，
\[
  \lim_{\substack{n\to\infty\\n\equiv6\ (\mathrm{mod}\ 12)}}
      \frac{f_{2/3}(n)}n=\frac14 .
\]
\end{theorem}
主定理中的 $T_0$ 一致于全部满足前提的 $E,A$。
因证明使用 Roth 的渐近结论，本文不提供数值上可用的显式阈值。
本稿尚未经过外部同行评审；附录中的部分形式化不替代主定理的纸面论证。

\section{记号及引用的增长定理}
对非空整数集 $B$ 与整数 $h\ge1$，记
\[
 hB=\{b_1+\cdots+b_h:b_i\in B\},
\]
其中\textbf{允许重复选取元素}。若 $0\in B$，则
$hB\subseteq(h+1)B$，且 $j(hB)=(jh)B$。
这与前述 $\sig(A)$ 必须严格区分。

\begin{lemma}[Lev 增量界；外部输入]\label{lem:lev}
设整数 $L\ge1$，有限整数集 $B\subseteq[0,L]$ 满足
$0,L\in B$ 与 $\gcdop(B)=1$，记 $N=|B|$。
则对每个整数 $j\ge2$，
\[
 |jB|\ge|(j-1)B|+\min\{L,\ j(N-2)+1\}.
 \tag{2.1}\label{eq:lev}
\]
\end{lemma}
这是 Lev~\cite{Lev1996} 的 Theorem~1。
其 Addendum~\cite{LevAddendum1997} 允许各个加数集合不同；
取所有加数集合均为 $B$ 时也得到同一不等式。
两文作者自存原文的定理与前提均已直接核对。
Matijasevick 的 Remark~3.2 提到的 Lev 文献
\cite{Lev1997} 则是另一篇关于和集区间的工作。

\begin{lemma}[有限步密度增长]\label{lem:growth}
设 $q\ge2$、$L\ge q$ 为整数，$B\subseteq[0,L]$ 为整数集，
$0,L\in B$、$\gcdop(B)=1$，且
\[
 |B|=N\ge \frac Lq+2.
\]
令 $r=q-1$。则
\[
 \boxed{|rB|\ge\frac{rL}{2}+q.}\tag{2.2}\label{eq:growth}
\]
\end{lemma}
\begin{proof}
若 $q=2$，则 $r=1$，结论正是 $N\ge L/2+2$。
若 $q\ge3$，对每个 $2\le j\le r$，有
\[
 j(N-2)+1\ge\frac{jL}{q}+1,
 \qquad
 \frac{jL}{q}+1\le L.
\]
第二式利用 $j\le q-1$ 与 $L\ge q$。
由引理~\ref{lem:lev}，
\[
 |jB|-|(j-1)B|\ge\frac{jL}{q}+1.
\]
从 $|B|\ge L/q+2$ 逐步相加，得到
\begin{align*}
 |rB|
 &\ge\frac Lq+2+\sum_{j=2}^{r}
                       \left(\frac{jL}{q}+1\right)\\
 &=\frac Lq\sum_{j=1}^{r}j+(r+1)
 =\frac{r(r+1)}{2q}L+(r+1)
 =\frac{rL}{2}+q .
\end{align*}
\end{proof}
特别地，$q=4$ 时 $r=3$，
$|B|\ge L/4+2$ 推出 $|3B|\ge3L/2+4$；
此处的 $3B$ 仍是允许重复的和集。

\section{中心区间覆盖与有界加法基}
\begin{lemma}[中心区间覆盖]\label{lem:central}
设 $D\ge1$ 是整数，$C\subseteq[0,D]$，
$0,D\in C$，且 $D\le2|C|-4$。
令
\[
 \gamma=2D+2-2|C|.
\]
则 $0\le\gamma\le D-2$，并且对每个整数 $j\ge2$，
\[
  [\gamma,jD-\gamma]\subseteq jC. \tag{3.1}\label{eq:central}
\]
该结论不需要 $\gcdop(C)=1$。
\end{lemma}
\begin{proof}
令 $N=|C|$。由于 $N\le D+1$，有 $\gamma\ge0$；
$D\le2N-4$ 恰等价于 $\gamma\le D-2$。
$[0,D]\setminus C$ 恰有 $\gamma/2$ 个元素。

先证二重和集覆盖。对
$z\in[\gamma,2D-\gamma]$，令
$U_z=\{u\in[0,D]:z-u\in[0,D]\}$。
则 $|U_z|=\min\{z,2D-z,D\}+1\ge\gamma+1$。
导致 $u\notin C$ 或 $z-u\notin C$ 的候选各不超过
$\gamma/2$ 个，故至少有一个 $u$ 同时满足
$u,z-u\in C$。于是
\[
 [\gamma,2D-\gamma]\subseteq2C. \tag{3.2}
\]

再令 $C'=C\setminus\{D\}\subseteq[0,D-1]$。
由于 $|C'|=N-1>D/2$，模 $D$ 的两份 $C'$
之和遍历所有剩余类；从而每个 $a\in[0,D-1]$
均可写为 $u+v=a$ 或 $u+v=a+D$，
其中 $u,v\in C'$。

给定 $j\ge2$ 与 $z\in[\gamma,jD-\gamma]$。
若 $z\le2D-\gamma$，由 (3.2) 得二项表示，
再补 $j-2$ 个零。
若 $z\ge(j-2)D+\gamma$，
则 $z-(j-2)D\in[\gamma,2D-\gamma]$，
可用二项表示并补 $j-2$ 个 $D$。
余下的中间情形满足
$2D-\gamma<z<(j-2)D+\gamma$。
写 $z=kD+a$，$0\le a<D$。
由 $\gamma\le D-2$ 得 $1\le k\le j-2$。
用上段的 $u,v\in C'$ 表示 $a$ 或 $a+D$；
前者补 $k$ 个 $D$，后者补 $k-1$ 个 $D$，
再用零补齐到恰好 $j$ 项。
故所有情形均有 $z\in jC$。
\end{proof}

\begin{lemma}[有界加法基]\label{lem:basis}
在引理~\ref{lem:growth} 的假设下，任给实数 $K>0$，
置
\[
 J=\lceil K\rceil+1,\qquad H=(q-1)J.
\]
对每个整数 $y$，若
$(q-1)L\le y\le KL$，则 $y\in HB$。
\end{lemma}
\begin{proof}
设 $r=q-1$、$C=rB$、$D=rL$。
显然 $C\subseteq[0,D]$ 且 $0,D\in C$。
由式~\eqref{eq:growth}，
\[
 |C|\ge D/2+q\ge D/2+2.
\]
故引理~\ref{lem:central} 适用，给出
$0\le\gamma\le D-2$ 和
$[\gamma,JD-\gamma]\subseteq JC$。
因为 $J\ge2$ 且 $r\ge1$，
\[
 y\ge D>\gamma,\qquad
 y\le KL\le\lceil K\rceil L
     \le(J-1)D\le JD-\gamma.
\]
于是 $y\in JC=(Jr)B=HB$。
注意得到的是\emph{恰好 $H$ 项}的允许重复表示；
零项填充是合法的。
\end{proof}

\begin{remark}[与 Lev 1997 区间定理的关系]\label{rem:lev1997}
引理~\ref{lem:basis} 也是 Lev~\cite{Lev1997}, Theorem~1$'$ 的直接推论。
为核对这一关系，令 $a=N-2$、$k=\lfloor(L-1)/a\rfloor$。
上述假设给出 $N\ge3$、$1\le a\le L-1$，
且 $1\le k\le q-1=r$；整除边界取该较大的整数选择亦符合原文条件。
该定理设
\[
 G=k\bigl(2L-2-(k+1)(N-2)\bigr),
\]
并对所有整数 $h\ge2k$ 给出 $[G,hL-G]\subseteq hB$。
由 $a\ge L/q$ 及
\[
 r-\frac{k(2q-k-1)}q=\frac{(q-k)(q-1-k)}q\ge0
\]
得 $G\le rL$。
另一方面 $H=r(\lceil K\rceil+1)\ge2r\ge2k$，
且 $HL-G\ge(H-r)L\ge KL$。
故 $[rL,KL]\subseteq HB$。
因此，本文使用的有界基窗口不是新的独立结构定理；
上面的增量式证明是这一既有结论的自足推导。
一般有界基窗口的线性因子也不能任意降低。
例如固定 $q\ge2$，取整数 $b\ge2$、$L=qb$ 及
$B=\{0,q,2q,\ldots,qb\}\cup\{qb-1\}$。
则 $|B|=L/q+2$、$\gcdop(B)=1$，
但 $y=(q-1)(L-2)-1$ 不属于任意重和集：
因 $y\equiv1\pmod q$，表示中 $qb-1$ 的份数至少为 $q-1$，
仅这些项的和就大于 $y$。
这个已有 Frobenius 型例族的代入说明一般窗口因子趋近 $q-1$；
它\emph{不}证明 $\theta_q=q-1$，因为后者的目标 $sT$ 另有整除结构。
\end{remark}

\section{由重复和到互异子集和：好核与删点}
对有限 $A\subseteq[1,E]$、$g\in A$，记
\[
 \nu_A(g)=
 \bigl|\{u\in A:u<g,\ 2g-u\in A\}\bigr|,\qquad
 C_m(A)=\{g\in A:\nu_A(g)<m\}.
\]
$r_3(E)$ 记 $[1,E]$ 中不含非平凡三项等差数列的
集合的最大基数。Roth 定理~\cite{Roth1953} 给出
$r_3(E)=o(E)$。

\begin{lemma}[重复项消去]\label{lem:distinct}
设整数 $m\ge1$，$G\subseteq A$，每个 $g\in G$
均满足 $\nu_A(g)\ge m$。
若 $t$ 可写成 $1\le h\le m$ 个（可重复）$G$ 元素之和，
则 $t\in\sig(A)$，且可用恰好 $h$ 个互异元素表示。
\end{lemma}
\begin{proof}
从给定多重表示开始；维持如下不变式：
有重复的值仅来自 $G$。
若 $g\in G$ 至少出现两次，删除两份 $g$。
其余表示中至多有 $h-2$ 个被占据的不同值。
以 $g$ 为中心的 $\nu_A(g)\ge m\ge h$ 个
无重合的二元配对 $\{u,2g-u\}$，
至少有一对与其余被占据的值不相交。
用这两个互异 $A$ 元素替代两份 $g$，
既不改变总和，也不改变总项数。
新添的两个元素均只出现一次，
令多重表示为 $\mu$，取自然数度量
$\delta(\mu)=h-|\operatorname{supp}\mu|$。
删除两份 $g$ 后支持集至多损失一个值，随后加入两个新值，
故支持集严格增大，$\delta$ 严格下降。
有限次操作后得到恰好 $h$ 项的互异表示；由 $h\ge1$ 可知其非空。
新添元素不必属于 $G$，
因为它们绝不会成为新的重复值。
\end{proof}

\begin{lemma}[坏中心规模]\label{lem:bad}
对整数 $m\ge1$ 及 $A\subseteq[1,E]$，
\[
 |C_m(A)|\le2\sqrt{2m}\,r_3(E). \tag{4.1}\label{eq:bad}
\]
\end{lemma}
\begin{proof}
记 $V=C_m(A)$、$v=|V|$。
在 $V$ 上建立以非平凡三项等差数列为边的三元超图；
按中心 $g$ 计数，边数至多 $(m-1)v$。
当 $m=1$，没有边，故 $v\le r_3(E)$，
从而 (4.1) 成立。
当 $m\ge2$，取
$p=(2(m-1))^{-1/2}\le1$，
以概率 $p$ 独立选取每个顶点。
所选顶点数减去完全落入的边数的期望至少为
\[
 pv-p^3(m-1)v=\frac{pv}{2}.
\]
因此可选到一个诱导子集，删除每条边至少一个顶点后，
仍有至少 $pv/2$ 个顶点且不含三项等差数列。
由 $pv/2\le r_3(E)$，
\[
 v\le2\sqrt{2(m-1)}\,r_3(E)
       \le2\sqrt{2m}\,r_3(E).\qedhere
\]
\end{proof}
重复项消去与坏中心界分别重述
Matijasevick~\cite{Matijasevick2026} 的 Lemma~2.2(a),(b)；
中心覆盖来自该文的 Lemma~2.3，好核及归一化也沿用其主证明框架。
这里给出完整证明，以明确重复和集与互异子集和的区别；
不将这些桥接引理作为本文的新增定理。

\section{线性范围逆定理的证明}
\begin{proof}[定理~\ref{thm:main} 的证明]
固定 $s\ge1,\eps>0$，令 $q=\snd(s)$、$r=q-1$。
\emph{先于} $T,E,A$ 固定
\[
 K=\frac{2s}{\eps},\quad
 J=\lceil K\rceil+1,\quad
 H=rJ,\quad m=H.
\]
由 $r_3(E)=o(E)$，选 $E_0=E_0(s,\eps)$，
使所有整数 $E\ge E_0$ 均满足
\[
  2\sqrt{2H}\,r_3(E)
       \le\frac{\eps r}{2s}E. \tag{5.1}
\]
再选整数 $T_0$，使
$\eps T_0\ge E_0$，并且 $\eps T_0/2\ge q$。
以下任取满足定理前提的 $T,E,A$。

\textbf{第一步：抽取好核。}
因为 $E\ge|A|\ge\eps T\ge E_0$，
令 $G=A\setminus C_H(A)$。
由引理~\ref{lem:bad} 与 (5.1)，
\[
 |A\setminus G|
 \le\frac{\eps r}{2s}E\le\frac{\eps T}{2}.
\]
因此
\[
 |G|\ge\frac Eq+\frac{\eps T}{2}>\frac Eq,
 \qquad \nu_A(g)\ge H\quad(g\in G). \tag{5.2}
\]

\textbf{第二步：处理公因子。}
令 $d=\gcdop(G)$，此时 $G\ne\varnothing$。
由于 $G$ 是区间 $[1,E]$ 内互异的正 $d$ 倍数，
$|G|\le E/d$；由 $|G|>E/q$ 得 $d<q$。
按 $q$ 的定义有 $d\mid s$，故
\[
   y=\frac{sT}{d}\in\N. \tag{5.3}
\]

\textbf{第三步：归一化及密度余量。}
置
\[
 B=\{0\}\cup\{g/d:g\in G\},\qquad
 L=\frac{\max G}{d}.
\]
则 $0,L\in B\subseteq[0,L]$、
$\gcdop(B)=1$、$|B|=|G|+1$，并且
\[
 L\ge|G|\ge\frac{\eps T}{2}\ge q,\qquad
 L\le E.
\]
须特别检查引理~\ref{lem:growth} 中的\textbf{$+2$}：
仅凭 $|G|>E/q$ 并不能推出它；
这里真正使用的是 (5.2) 中保留的余量
$\eps T/2\ge q$。故
\[
 |B|=|G|+1
 \ge \frac Eq+\frac{\eps T}{2}+1
 \ge \frac Lq+q+1
 \ge \frac Lq+2. \tag{5.4}
\]
另一方面，由 $\max G\le E$、$\max G\ge|G|$，
\[
 r\le\frac{sT}{E}
 \le\frac{y}{L}
 =\frac{sT}{\max G}
 \le\frac{sT}{|G|}
 \le\frac{2s}{\eps}=K. \tag{5.5}
\]
引理~\ref{lem:basis} 用 (5.4)--(5.5) 给出
$y\in HB$。

\textbf{第四步：回到互异子集和。}
先将 $y$ 的恰好 $H$ 项 $B$ 表示删去零项，
再乘以 $d$，得到 $sT$ 由至多 $H$ 个
（允许重复的）$G$ 元素表示。因为 $sT>0$，删零后项数满足 $1\le h\le H$。
结合 (5.2) 与引理~\ref{lem:distinct}，
得到 $sT$ 的至多 $H$ 个互异 $A$ 元素表示。
所有依赖于 $T,E,A$ 的选择都在固定 $H,E_0,T_0$
之后完成，故量词一致。证毕。
\end{proof}

\begin{corollary}[范围常数]\label{cor:theta}
在 Matijasevick~\cite{Matijasevick2026} 对 $\theta_q$ 的定义下，
定理~\ref{thm:main} 蕴含
\[
 \theta_q\le q-1.
\]
因此对可出现的素数幂 $q=\snd(s)\ge4$，
\[
 \boxed{\max\!\left\{2,q-2+\frac1q\right\}
        \le\theta_q\le q-1}.
\]
特别地 $\theta_4\in[9/4,3]$。
而 $q=2,3$ 与原文已知的 $\theta_2=1,\theta_3=2$ 一致。
\end{corollary}
\begin{proof}
$\theta_q$ 是所有使定理的范围条件
$\theta E\le sT$ 成立的参数 $\theta$ 的下确界。
定理~\ref{thm:main} 的范围因子为 $q-1$，
并且其表示长度界不大于原 Theorem~1.2 的
$2^{q-2}(\lceil2s/\eps\rceil+1)$，
因此确实满足 Proposition~3.1 所定义参数的全部结论，
而不只是无长度限制的存在性结论。
故 $\theta=q-1$ 可用；下界沿用原文的已知构造。
又 $q-1\le2^{q-2}$，因此新参数不比原范围更弱。
\end{proof}

\section{\texorpdfstring{$q=4$ 与 $c=2/3$}{q=4 与 c=2/3} 的指定同余类极限}
\begin{proof}[定理~\ref{thm:limit} 的证明]
将 $n$ 写成 $n=6T$，其中 $T$ 为正奇数。
则 $n\equiv6\pmod{12}$，且
$\lfloor2n/3\rfloor=4T$。
此时 $\snd(6)=4$；对每个奇数 $T$，
$6T\equiv2\pmod4$。

\textbf{上界。} 设 $A\subseteq[1,4T]$
满足 $6T\notin\sig(A)$。
分解
\[
 A_0=A\cap[1,2T-1],\qquad
 A_1=A\cap[2T,4T].
\]
给定任意 $\eps>0$，在定理~\ref{thm:main} 中取
$s=6,q=4,r=3,E=2T-1$。
由于 $3(2T-1)\le6T$，
对充分大的奇数 $T$，若
$|A_0|\ge(2T-1)/4+\eps T$，则
$6T\in\sig(A_0)$，矛盾。故
\[
 |A_0|<\frac{2T-1}{4}+\eps T
        <\left(\frac12+\eps\right)T. \tag{6.1}
\]
另一方面，沿用 \cite[Lemma~5.1]{Matijasevick2026} 的配对，
$[2T,4T]$ 被 $T$ 对互补元素
$\{x,6T-x\}$（$2T\le x\le3T-1$）
以及单点 $3T$ 分割。
$A_1$ 在每对中至多选一个，因此
\[
 |A_1|\le T+1. \tag{6.2}
\]
由 (6.1)--(6.2)，
\[
 F(4T,6T)\le\left(\frac32+\eps\right)T+1.
\]
先令 $T\to\infty$（沿奇数），再令 $\eps\downarrow0$，
得到
\[
 \limsup_{\substack{T\to\infty\\T\text{ 为奇数}}}
 \frac{F(4T,6T)}{6T}\le\frac14. \tag{6.3}
\]

\textbf{下界。} 取
\[
 Q=\{4,8,\ldots,4T\},\quad
 U=\{x\in[3T,4T]:x\ \text{为奇数}\},\quad
 A=Q\cup U.
\]
$Q$ 的任意子集和模 $4$ 为 $0$，
而 $6T\equiv2\pmod4$；
若子集包含恰好一个 $U$ 元素，其和为奇数，
不可能等于偶数 $6T$；
若包含至少两个互异 $U$ 元素，其和严格大于
$2\cdot3T=6T$。
故 $6T\notin\sig(A)$。
又因 $T$ 为奇数，
\[
 |Q|=T,\qquad |U|=\frac{T+1}{2},
 \qquad |A|=\frac{3T+1}{2}.
\]
于是
\[
 \liminf_{\substack{T\to\infty\\T\text{ 为奇数}}}
 \frac{F(4T,6T)}{6T}\ge\frac14. \tag{6.4}
\]
由 (6.3)--(6.4) 得所求极限。
此下界构造是原文 Proposition~5.3(b)
的模数构造的 $q=4$ 特例，并非新构造。
\end{proof}


\section{结语与适用边界}
本证明的变化集中在一个有限步增长估计：
Lev 增量界使 $(q-1)B$ 跨过半密度，再利用中心区间覆盖给出有界项数表示。
好中心及消重框架把重复表示转回原集合中的互异表示。
当 $q\ge4$ 时，$q-1<2^{q-2}$，范围和长度均严格改进；
$q=2,3$ 时与原范围一致。
本文只详证 $q=4,c=2/3$ 的同余类应用，
不把整个中间参数区间的极限公式列为已证结论。
这是一份供外部审阅的未发表文稿；
有关贡献与既有结果的对照见随附的文献比较说明，
有界检索不作为新颖性或优先权的证明。

\appendix
\section{可复现核验及形式化边界}\label{sec:verify}
附带材料包含仅用 Python 标准库的有限检查脚本、
去重桥与算术步骤的 Lean 源码、运行说明及文件哈希。
有限检查覆盖小集合的和集增长、中心区间、有界基、消重、坏中心与应用构造，
包含去掉 gcd 条件和削弱密度余量的负控制。
独立构造检查另外覆盖空表示、消重支持增长及模 $D$ 填充的中间分支。
这些结果辅助检错，不证明本文的渐近命题，也不替代所引定理。

Lean 文件 \texttt{EStarDistinct.lean} 对有限自然数集和多重集证明：
若 $G\subseteq A$、每个中心有至少 $m$ 个配对，
任意至多 $m$ 项的 $G$ 重复和可转为同和、同项数的互异 $A$ 子集。
非空推论显式要求正项数。
文件中的五条相关声明及原有九条算术声明已由固定工具链检查，
公理报告仅含 \texttt{propext, Classical.choice, Quot.sound}；
本次独立重实现的组合桥模块另经 \texttt{leanchecker} 重放。
原预印本亦报告该组合桥的形式化；这里不主张首次形式化。
这不包含 Lev、Roth、坏中心概率估计、统一阈值及最终定理的全部形式化，
不应被解释为主定理的端到端机器证明。
复现包中精确记录有限范围、版本、退出码及未覆盖项。

\begin{thebibliography}{99}
\bibitem{Matijasevick2026}
M.~Matijasevick,
\emph{Forbidden subset sums in a linear range: an inverse theorem and the Erd\H{o}s--Graham irregularity problem},
arXiv:2611.17205v1, 2026.
\url{https://arxiv.org/abs/2611.17205v1}.

\bibitem{Lev1996}
V.~F.~Lev,
\emph{Structure theorem for multiple addition and the Frobenius problem},
Journal of Number Theory \textbf{58} (1996), 79--88.
\href{https://doi.org/10.1006/jnth.1996.0065}{doi:10.1006/jnth.1996.0065}.
作者自存版：\url{https://math.haifa.ac.il/seva/Papers/multiad.dvi}.

\bibitem{LevAddendum1997}
V.~F.~Lev,
\emph{Addendum to ``Structure Theorem for Multiple Addition''},
Journal of Number Theory \textbf{65} (1997), 96--100.
\href{https://doi.org/10.1006/jnth.1997.2128}{doi:10.1006/jnth.1997.2128}.
作者自存版：\url{https://math.haifa.ac.il/seva/Papers/adndum.dvi}.

\bibitem{Lev1997}
V.~F.~Lev,
\emph{Optimal representations by sumsets and subset sums},
Journal of Number Theory \textbf{62} (1997), 127--143.
\href{https://doi.org/10.1006/jnth.1997.2012}{doi:10.1006/jnth.1997.2012}.

\bibitem{Roth1953}
K.~F.~Roth,
\emph{On certain sets of integers},
Journal of the London Mathematical Society \textbf{s1-28} (1953), no.~1, 104--109.
\href{https://doi.org/10.1112/jlms/s1-28.1.104}{doi:10.1112/jlms/s1-28.1.104}.
\end{thebibliography}
\end{document}
